% !TEX root = ../main.tex
\begin{ouverture}[doppler]{23}{Effet Doppler}{Partie 4 · Ondes et signaux}
{Émetteur en mouvement\enspace·\enspace Décalage Doppler\enspace·\enspace Mesure de vitesses\enspace·\enspace Décalage vers le rouge}
\objectif{Décrire et interpréter qualitativement les observations correspondant à une manifestation de l'effet Doppler.}
\objectif{Établir l'expression du décalage Doppler pour un observateur fixe et un émetteur mobile, à une dimension.}
\objectif{Exploiter l'expression du décalage Doppler pour des ondes acoustiques et électromagnétiques.}
\objectif{Déterminer une vitesse à partir d'un décalage Doppler mesuré.}
\objectif{Calculer une vitesse et une fréquence émise à partir de deux fréquences mesurées, en Python.}
\prerequis{Période, fréquence, longueur d'onde, célérité : $\lambda=v\,T$ (classe de première).}
\prerequis{Ondes électromagnétiques, spectres de raies (classes de seconde et de première).}
\prerequis{Mouvement rectiligne uniforme (chapitre~13).}
\end{ouverture}
\essentiel
\essSection{1}{Le phénomène}
\begin{definition}[Effet Doppler]
Lorsqu'un émetteur d'ondes se déplace par rapport à un récepteur, la fréquence reçue $f_\mathrm{R}$ diffère de la fréquence
émise $f_\mathrm{E}$ : elle est plus élevée quand l'émetteur se rapproche, plus basse quand il s'éloigne. C'est l'\textbf{effet
Doppler}. Le son de la sirène d'une ambulance paraît plus aigu à son approche, plus grave après son passage. La
différence $\Delta f=f_\mathrm{R}-f_\mathrm{E}$ est le \textbf{décalage Doppler}.
\end{definition}
\schema{23-fronts}
\essSection{2}{L'expression du décalage Doppler}
\begin{savoirfaire}[Établir l'expression de la fréquence reçue]
Un émetteur se rapproche à la vitesse $v_\mathrm{E}$ d'un récepteur fixe, sur une même droite ; l'onde se propage à la
célérité $v$. L'émetteur émet un maximum à la date 0, puis le suivant à la date $T_\mathrm{E}$. Entre-temps, le premier
maximum a parcouru $v\,T_\mathrm{E}$ et l'émetteur $v_\mathrm{E}\,T_\mathrm{E}$ : les deux maximums sont séparés de
\[\lambda_\mathrm{R}=(v-v_\mathrm{E})\,T_\mathrm{E}\qquad\text{d'où}\qquad f_\mathrm{R}=\frac{v}{\lambda_\mathrm{R}}=f_\mathrm{E}\,\frac{v}{v-v_\mathrm{E}}.\]
Si l'émetteur s'éloigne, on remplace $v_\mathrm{E}$ par $-v_\mathrm{E}$ : $f_\mathrm{R}=f_\mathrm{E}\,\dfrac{v}{v+v_\mathrm{E}}$.
\end{savoirfaire}
\schema{23-demonstration}
\begin{loi}[Décalage Doppler]
Pour un émetteur qui se rapproche d'un récepteur fixe,
\[\Delta f=f_\mathrm{R}-f_\mathrm{E}=f_\mathrm{E}\,\frac{v_\mathrm{E}}{v-v_\mathrm{E}}\approx f_\mathrm{E}\,\frac{v_\mathrm{E}}{v}\quad\text{si }v_\mathrm{E}\ll v ;\]
le décalage est négatif si l'émetteur s'éloigne. Le décalage relatif $\dfrac{|\Delta f|}{f_\mathrm{E}}$ est voisin de $\dfrac{v_\mathrm{E}}{v}$ :
mesurer un décalage permet de mesurer une vitesse.
\end{loi}
\begin{exemple}
Une voiture roulant à \qty{90}{km.h^{-1}} (\qty{25}{m.s^{-1}}) klaxonne à \qty{900}{Hz}. Le son, de célérité \qty{340}{m.s^{-1}}, est
perçu à $900\times\dfrac{340}{315}=\qty{971,4}{Hz}$ à son approche et à $900\times\dfrac{340}{365}=\qty{838,4}{Hz}$ quand elle s'éloigne.
L'approximation donne $\Delta f\approx\qty{66}{Hz}$ au lieu de \qty{71,4}{Hz} : elle n'est valable que si $v_\mathrm{E}\ll v$.
\end{exemple}
\essSection{3}{Des vitesses mesurées par effet Doppler}
\begin{propriete}[Applications]
\begin{itemize}
\item Les \textbf{radars routiers} émettent une onde électromagnétique qui se réfléchit sur le véhicule ; la réflexion
sur une cible mobile double le décalage : $\Delta f\approx2\,f\,\dfrac{v_\mathrm{E}}{c}$ pour une cible qui se déplace dans l'axe.
\item L'\textbf{échographie Doppler} mesure la vitesse du sang dans les vaisseaux par réflexion d'ultrasons.
\item En \textbf{astrophysique}, les raies du spectre d'une étoile ou d'une galaxie qui s'éloigne sont décalées vers les
grandes longueurs d'onde (« vers le rouge ») : $\dfrac{\lambda_\mathrm{R}-\lambda_\mathrm{E}}{\lambda_\mathrm{E}}\approx\dfrac{v_\mathrm{E}}{c}$. Ce décalage a
révélé l'expansion de l'Univers et permet de détecter des planètes autour d'autres étoiles.
\end{itemize}
\end{propriete}
\schema{23-spectre}
\begin{codecote}[0.64]
\begin{pythoncol}
v = 340.0                         # m/s
f_app, f_eloi = 971.4, 838.4      # Hz
vE = v * (f_app - f_eloi) / (f_app + f_eloi)
fE = f_app * (v - vE) / v
print(round(vE, 1), round(3.6 * vE), round(fE))
\end{pythoncol}
\begin{consolecol}
25.0 90 900
\end{consolecol}
\cote
\fontsize{9.5}{12.5}\selectfont
\capacite{Capacité numérique}\par\smallskip
À partir des fréquences mesurées à l'approche et à l'éloignement, le programme calcule la vitesse de l'émetteur, en
\unit{m.s^{-1}} et en \unit{km.h^{-1}}, puis la fréquence émise : on retrouve la voiture de l'exemple, sans connaître à
l'avance la fréquence de son klaxon.
\end{codecote}
\begin{attention}
Les relations $f_\mathrm{R}=f_\mathrm{E}\,\dfrac{v}{v\mp v_\mathrm{E}}$ supposent un récepteur fixe dans le milieu de propagation et un
émetteur qui se déplace sur la droite qui les joint. Pour une onde réfléchie par une cible mobile, comme dans un radar,
le décalage est doublé.
\end{attention}
\begin{experience}[Mesurer une vitesse par effet Doppler]
Un émetteur d'ultrasons de \qty{40}{kHz}, fixé sur un chariot, se rapproche puis s'éloigne d'un récepteur relié à une
interface d'acquisition. L'analyse spectrale du signal reçu donne la fréquence pendant chaque phase ; on en déduit la
vitesse du chariot, que l'on compare à celle obtenue par chronophotographie. Une application de smartphone permet de
même d'étudier le son d'une voiture qui passe.
\end{experience}
\begin{rediger}
Préciser l'émetteur, le récepteur et le sens du mouvement : « l'émetteur se rapproche, donc $f_\mathrm{R}>f_\mathrm{E}$ ».
Écrire la relation exacte avant d'utiliser l'approximation $v_\mathrm{E}\ll v$, et justifier cette dernière par un ordre de
grandeur.
\end{rediger}
\exercices

\rubrique{Découverte}
\begin{blocLivre}{ex-23-1}
\begin{exercice}[1]{La sirène des pompiers}
La sirène d'un camion de pompiers émet alternativement deux notes de fréquences \qty{435}{Hz} et \qty{488}{Hz}. Le camion roule
à \qty{72}{km.h^{-1}}.
\begin{donnees}
Célérité du son dans l'air : $v=\qty{340}{m.s^{-1}}$.
\end{donnees}
\begin{questions}
\item Décrire ce qu'entend un piéton immobile pendant que le camion se rapproche, puis s'éloigne.
\item Calculer la fréquence perçue pour la note de \qty{435}{Hz} à l'approche, puis à l'éloignement.
\item Calculer la fréquence perçue à l'approche pour la note de \qty{488}{Hz}.
\item Les pompiers, assis dans le camion, entendent-ils un effet Doppler ? Justifier.
\end{questions}
\end{exercice}
\end{blocLivre}
\begin{blocLivre}{ex-23-2}
\begin{exercice}[1]{Un émetteur qui s'éloigne}
Un émetteur s'éloigne à la vitesse $v_\mathrm{E}$ d'un récepteur fixe ; l'onde se propage à la célérité $v$.
\begin{questions}
\item L'émetteur émet un maximum à la date 0 et le suivant à la date $T_\mathrm{E}$. Quelle distance l'émetteur a-t-il parcourue
entre ces deux dates ? Dans quel sens ?
\item Montrer que deux maximums successifs reçus sont séparés de $\lambda_\mathrm{R}=(v+v_\mathrm{E})\,T_\mathrm{E}$.
\item En déduire l'expression de la fréquence reçue en fonction de $f_\mathrm{E}$, $v$ et $v_\mathrm{E}$.
\item La fréquence reçue est-elle plus grande ou plus petite que la fréquence émise ? Retrouver ce résultat sur le
schéma des fronts d'onde de l'essentiel.
\end{questions}
\schema{23-fronts}
\end{exercice}
\end{blocLivre}
\begin{blocLivre}{ex-23-3}
\begin{exercice}[1]{Le sifflet du train}
Un train roulant à \qty{144}{km.h^{-1}} siffle en traversant une gare ; le sifflet émet un son de \qty{500}{Hz}.
\begin{donnees}
$v=\qty{340}{m.s^{-1}}$.
\end{donnees}
\begin{questions}
\item Convertir la vitesse du train en \unit{m.s^{-1}}.
\item Calculer la fréquence perçue par un voyageur sur le quai à l'approche du train, puis après son passage.
\item Calculer le rapport de ces deux fréquences. Le changement de hauteur est-il perceptible ?
\item Le sifflet est-il entendu à la même fréquence par le conducteur ? Justifier.
\end{questions}
\end{exercice}
\end{blocLivre}
\begin{blocLivre}{ex-23-4}
\begin{exercice}[1]{Quand l'approximation est-elle valable ?}
Un émetteur de fréquence \qty{1000}{Hz} se rapproche d'un récepteur ; le son se propage à \qty{340}{m.s^{-1}}.
\begin{questions}
\item Calculer le décalage Doppler exact pour $v_\mathrm{E}=\qty{10}{m.s^{-1}}$, \qty{50}{m.s^{-1}} et \qty{100}{m.s^{-1}}.
\item Calculer pour les mêmes vitesses le décalage approché $f_\mathrm{E}\,v_\mathrm{E}/v$.
\item Calculer l'écart relatif entre les deux valeurs pour \qty{10}{m.s^{-1}}. Conclure sur la validité de l'approximation.
\item Pourquoi l'approximation est-elle toujours excellente pour les ondes lumineuses et les véhicules terrestres ?
\end{questions}
\end{exercice}
\end{blocLivre}
\begin{blocLivre}{ex-23-5}
\begin{exercice}[1]{Une étoile qui s'éloigne}
Au laboratoire, la raie H$\alpha$ de l'hydrogène a pour longueur d'onde \qty{656,28}{nm}. Dans le spectre d'une étoile, elle est
mesurée à \qty{656,50}{nm}.
\begin{donnees}
$c=\qty{3,00e8}{m.s^{-1}}$ ; $\dfrac{\lambda_\mathrm{R}-\lambda_\mathrm{E}}{\lambda_\mathrm{E}}\approx\dfrac{v_\mathrm{E}}{c}$.
\end{donnees}
\begin{questions}
\item La raie est-elle décalée vers le rouge ou vers le bleu ? L'étoile se rapproche-t-elle ou s'éloigne-t-elle ?
\item Calculer la vitesse de l'étoile par rapport à la Terre, en \unit{km.s^{-1}}.
\item Pourquoi doit-on mesurer la longueur d'onde avec une très grande précision ?
\item Quelle composante de la vitesse de l'étoile cette méthode mesure-t-elle ?
\end{questions}
\end{exercice}
\end{blocLivre}
\rubrique{Entraînement}
\begin{blocLivre}{ex-23-6}
\begin{exercice}[2]{Le radar routier}
Un radar fixe émet une onde électromagnétique de fréquence $f=\qty{24,125}{GHz}$ en direction des véhicules qui
arrivent ; son faisceau fait un angle $\alpha=25^\circ$ avec la route. Pour une onde réfléchie par un véhicule de vitesse $v_\mathrm{E}$,
le décalage vaut $\Delta f=\dfrac{2\,f\,v_\mathrm{E}\cos\alpha}{c}$. Pour une voiture, il mesure $\Delta f=\qty{3,95}{kHz}$.
\begin{donnees}
$c=\qty{3,00e8}{m.s^{-1}}$. Pour une vitesse inférieure à \qty{100}{km.h^{-1}}, on retient la vitesse mesurée diminuée de
\qty{5}{km.h^{-1}}.
\end{donnees}
\begin{questions}
\item Pourquoi le décalage fait-il intervenir le facteur 2 ? et le facteur $\cos\alpha$ ?
\item Calculer la vitesse de la voiture, en \unit{m.s^{-1}} puis en \unit{km.h^{-1}}.
\item La route est limitée à \qty{90}{km.h^{-1}}. Le conducteur est-il en infraction ?
\item Quel décalage le radar mesurerait-il pour une voiture roulant exactement à \qty{90}{km.h^{-1}} ?
\end{questions}
\end{exercice}
\end{blocLivre}
\begin{blocLivre}{ex-23-7}
\begin{exercice}[2]{La vitesse d'une voiture au bord de la route}
On reprend le programme de l'essentiel. Un élève enregistre le bruit d'une voiture qui passe : le son le plus intense
a pour fréquence \qty{1062}{Hz} à l'approche et \qty{878}{Hz} à l'éloignement.
\begin{questions}
\item \capacite{Capacité numérique}Montrer que les deux fréquences mesurées vérifient $f_\text{app}\,(v-v_\mathrm{E})=f_\text{éloi}\,(v+v_\mathrm{E})$,
puis justifier la ligne de calcul de \code{vE}.
\item Justifier la ligne de calcul de \code{fE}.
\item Quelles valeurs faut-il modifier dans le programme ? Calculer la vitesse de la voiture et la fréquence émise.
\item L'élève était placé à quelques mètres de la route. Pourquoi les fréquences mesurées juste au moment du passage ne
doivent-elles pas être utilisées ?
\end{questions}
\end{exercice}
\end{blocLivre}
\begin{blocLivre}{ex-23-8}
\begin{exercice}[2]{L'échographie Doppler}
Une sonde émet des ultrasons de \qty{4,0}{MHz} dans une artère. L'onde réfléchie par les globules rouges présente un
décalage $\Delta f=\dfrac{2\,f\,v\cos\theta}{c_\mathrm{t}}$, où $v$ est la vitesse du sang, $\theta=60^\circ$ l'angle entre le faisceau et
l'artère, et $c_\mathrm{t}=\qty{1540}{m.s^{-1}}$ la célérité des ultrasons dans les tissus.
\begin{questions}
\item L'appareil mesure $\Delta f=\qty{1,3}{kHz}$. Calculer la vitesse du sang.
\item Ce décalage est-il audible ? Les échographes le diffusent souvent par un haut-parleur : pourquoi est-ce utile ?
\item Quel décalage correspondrait à une vitesse de \qty{1,0}{m.s^{-1}}, signe d'un rétrécissement de l'artère ?
\item Pourquoi ne faut-il pas placer la sonde perpendiculairement à l'artère ?
\end{questions}
\end{exercice}
\end{blocLivre}
\begin{blocLivre}{ex-23-9}
\begin{exercice}[2]{Vers le mur du son}
Un avion émet un son de fréquence \qty{1000}{Hz} et vole vers un observateur à la vitesse $v_\mathrm{E}$. Le son se propage à
\qty{340}{m.s^{-1}}.
\begin{questions}
\item Calculer la longueur d'onde du son à l'arrêt, puis devant l'avion lorsque $v_\mathrm{E}=0{,}9\,v$.
\item Quelle fréquence l'observateur perçoit-il alors ?
\item Que devient $\lambda_\mathrm{R}$ lorsque $v_\mathrm{E}$ tend vers $v$ ? Décrire ce qui se passe pour les fronts d'onde.
\item Expliquer le « bang » entendu au passage d'un avion supersonique.
\end{questions}
\end{exercice}
\end{blocLivre}
\begin{blocLivre}{ex-23-10}
\begin{exercice}[2]{Une étoile double}
Dans le spectre d'une étoile appartenant à un système double, la raie H$\alpha$ oscille périodiquement entre \qty{656,10}{nm} et
\qty{656,46}{nm}. Au laboratoire, elle vaut \qty{656,28}{nm} ; le système, dans son ensemble, ne s'éloigne pas de la Terre.
\begin{donnees}
$c=\qty{3,00e8}{m.s^{-1}}$.
\end{donnees}
\begin{questions}
\item Interpréter l'oscillation de la raie.
\item Quand l'étoile se rapproche-t-elle de la Terre ? Quand s'en éloigne-t-elle ?
\item Calculer la valeur maximale de la vitesse de l'étoile dans la direction de la Terre.
\item Pourquoi ne détecte-t-on aucun décalage si l'orbite est vue de face ?
\end{questions}
\end{exercice}
\end{blocLivre}
\begin{blocLivre}{ex-23-11}
\begin{exercice}[2]{Le son d'une moto}
Un micro placé au bord d'une route enregistre le son d'une moto qui passe à vitesse constante.
\schema{23-ex11}
\begin{donnees}
$v=\qty{340}{m.s^{-1}}$.
\end{donnees}
\begin{questions}
\item Relever la fréquence perçue à l'approche et à l'éloignement de la moto.
\item Établir l'expression de la vitesse de la moto en fonction de ces deux fréquences, puis la calculer, en
\unit{m.s^{-1}} et en \unit{km.h^{-1}}.
\item Calculer la fréquence du son émis par la moto.
\item À quel moment la moto est-elle passée devant le micro ? Justifier.
\end{questions}
\end{exercice}
\end{blocLivre}
\rubrique{Objectif bac}
\begin{blocLivre}{ex-23-12}
\begin{exercice}[3]{L'expansion de l'Univers}
\begin{docu}{1}{La découverte de Hubble}
En 1929, Edwin Hubble constate que les galaxies lointaines s'éloignent de nous à une vitesse proportionnelle à leur
distance : $v=H_0\,d$. On détermine leur vitesse par le décalage vers le rouge de leurs raies.
\end{docu}
\begin{docu}{2}{Trois galaxies}
\begin{center}
\begin{tabular}{lccc}
\toprule
Galaxie & A & B & C \\
Distance (Mpc) & 20 & 50 & 100 \\
Raie H$\alpha$ mesurée (nm) & 659,4 & 663,9 & 671,7 \\
\bottomrule
\end{tabular}
\end{center}
Au laboratoire, la raie H$\alpha$ vaut \qty{656,3}{nm}. Un mégaparsec vaut \qty{3,086e19}{km}.
\end{docu}
\begin{donnees}
$c=\qty{3,00e5}{km.s^{-1}}$ ; une année vaut \qty{3,156e7}{s}.
\end{donnees}
\begin{questions}
\item Calculer la vitesse d'éloignement de chaque galaxie.
\item \initiative Les mesures vérifient-elles la loi de Hubble ? Estimer la valeur de $H_0$, en \unit{km.s^{-1}} par mégaparsec.
\item L'inverse $1/H_0$ donne un ordre de grandeur de l'âge de l'Univers. Le calculer en années, pour $H_0=\qty{70}{km.s^{-1}}$
par mégaparsec.
\item Une galaxie lointaine présente un décalage de la raie H$\alpha$ de \qty{70}{nm}. L'approximation utilisée reste-t-elle
valable ?
\end{questions}
\end{exercice}
\end{blocLivre}
\begin{blocLivre}{ex-23-13}
\begin{exercice}[3]{À la recherche des exoplanètes}
\begin{docu}{1}{Une étoile qui oscille}
Une planète et son étoile tournent autour de leur centre de masse commun : l'étoile se rapproche et s'éloigne
périodiquement de la Terre. En 1995, Michel Mayor et Didier Queloz détectent ainsi la première planète autour d'une
étoile semblable au Soleil, 51~Pegasi~b : la vitesse de l'étoile dans la direction de la Terre oscille avec une amplitude
de \qty{56}{m.s^{-1}} et une période de 4,23~jours.
\end{docu}
\begin{docu}{2}{Un spectrographe de précision}
Le spectrographe utilisé mesurait la position des raies d'absorption de l'étoile avec une précision relative de l'ordre
de $10^{-8}$.
\end{docu}
\begin{donnees}
$c=\qty{3,00e8}{m.s^{-1}}$.
\end{donnees}
\begin{questions}
\item Pourquoi l'étoile se déplace-t-elle ? Comment ce mouvement se traduit-il dans son spectre ?
\item Calculer le décalage relatif maximal des raies, puis le décalage maximal d'une raie de \qty{500}{nm}.
\item \initiative Le spectrographe permettait-il de détecter ce décalage ? Expliquer pourquoi cette découverte a demandé
une technologie exceptionnelle.
\item Jupiter fait osciller le Soleil à la vitesse maximale de \qty{12,5}{m.s^{-1}}. Une planète comme Jupiter serait-elle plus
facile ou plus difficile à détecter que 51~Pegasi~b ?
\end{questions}
\end{exercice}
\end{blocLivre}
\begin{blocLivre}{ex-23-14}
\begin{exercice}[3]{Écouter un satellite}
\begin{docu}{1}{Un satellite en orbite basse}
Un petit satellite universitaire émet une balise radio de fréquence \qty{437}{MHz}. Lors de son passage au-dessus d'une
station au sol, la composante de sa vitesse dans la direction de la station varie de \qty{+7}{km.s^{-1}} (il s'approche) à
\qty{-7}{km.s^{-1}} (il s'éloigne).
\end{docu}
\begin{docu}{2}{Le récepteur}
La bande passante du récepteur de la station n'est que de quelques kilohertz : un opérateur ou un logiciel doit ajuster
en permanence la fréquence d'écoute.
\end{docu}
\begin{donnees}
$c=\qty{3,00e8}{m.s^{-1}}$.
\end{donnees}
\begin{questions}
\item Calculer le décalage Doppler maximal, puis l'excursion totale de fréquence pendant un passage.
\item Pourquoi faut-il ajuster la fréquence d'écoute ?
\item \initiative Le satellite émet aussi des données à \qty{2,4}{GHz}. Le problème est-il plus ou moins gênant à cette fréquence ?
\item Au moment où le satellite passe au plus près de la station, quel est le décalage ? Justifier.
\end{questions}
\end{exercice}
\end{blocLivre}
\begin{blocLivre}{ex-23-15}
\begin{exercice}[3]{Le radar météorologique}
\begin{docu}{1}{Mesurer le vent dans les nuages}
Un radar météorologique émet une onde de \qty{5,6}{GHz}. Les gouttes de pluie emportées par le vent la réfléchissent ; le
décalage mesuré vaut $\Delta f=\dfrac{2\,f\,v_\mathrm{r}}{c}$, où $v_\mathrm{r}$ est la vitesse des gouttes dans la direction du radar,
comptée positivement quand elles se rapprochent.
\end{docu}
\begin{docu}{2}{Une mesure}
Dans la direction nord, à \qty{40}{km} du radar, le décalage vaut \qty{+560}{Hz} ; dans la direction sud, à la même distance, il
vaut \qty{-560}{Hz}.
\end{docu}
\begin{donnees}
$c=\qty{3,00e8}{m.s^{-1}}$.
\end{donnees}
\begin{questions}
\item Calculer la vitesse des gouttes dans la direction nord, en \unit{m.s^{-1}} et en \unit{km.h^{-1}}.
\item \initiative En exploitant les deux mesures, déterminer la direction et le sens du vent.
\item Quel décalage mesurerait-on pour une pluie poussée vers le radar à \qty{8}{m.s^{-1}} ?
\item Pourquoi le radar ne peut-il pas mesurer un vent soufflant perpendiculairement à son faisceau ?
\end{questions}
\end{exercice}
\end{blocLivre}
